% ===================================================================
%  ТАБЛИЦЯ № 7 — Село взимку. — Ферма навесні.
%  Traduction 2026 d'apres texto/io, controlee sur texto/fr, texto/en
%  et texto/es. Consigne : voir texto/uk/10-tablycia-01.tex.
%
%  OUVERTURE DE LA DEUXIEME SERIE. « ДРУГА СЕРІЯ » : le groupe
%  ORDINALO recoit « друга » et le groupe SERIO « серія », et
%  l'apparat se range de lui-meme sous le role serio.
%
%  LE LEXIQUE DE LA FERME EST LE PLUS UKRAINIEN DE TOUS, et le plus
%  eloigne du russe : le charron est un « стельмах », le bourrelier un
%  « лимар », le savetier un « латальник », le puits une « криниця »,
%  la grange une « клуня », le troupeau une « отара », la houlette une
%  « ґирлиґа », le journalier un « наймит ». Huit mots du meme cour de
%  ferme, huit mots que le russe ne dit pas ainsi.
% ===================================================================

%%K t07-noto-1 noto
\VUnotes{143.2mm}{%
(*) Подробиці обіду див. у таблицях 6 і 13.%
}

%%K t07-apar-1 apar
\VUsaut{4.0mm}
\VUcentre{13.2pt}{90}{ДРУГА СЕРІЯ}
\VUsaut{3.0mm}
\VUfilet{16mm}
\VUsaut{5.6mm}
\VUtitre{15.0pt}{60}{ТАБЛИЦЯ № 7}
\VUsaut{3.0mm}

%%K t07-tit-1 sub
\VUcentre{12.0pt}{0}{\textit{Перша сцена.}}
\VUsaut{3.0mm}

%%K t07-tit-2 sub
\VUpk{10.2pt}{40}{Село взимку.}
\VUsaut{4.0mm}

%%K t07-c1-01-1 p
1. --- На новорічні свята я поїхав з батьком до Новосілок, сільця, де в нього
були справи.

%%K t07-c1-02-1 p
2. --- Ми поїхали \VUgras{диліжансом} \textsuperscript{(11)}, який ходить між
містом і цим селом; але що було дуже холодно, подорож у такому незручному
возі вийшла зовсім не приємна.

%%K t07-c1-03-1 p
3. --- Тож ми щиро зраділи, коли \VUgras{кучер} спинив свій диліжанс на площі,
перед \VUgras{корчмою} \textsuperscript{(2)} \textit{«Білий Ведмідь»}.
\VUgras{Корчмар} \textsuperscript{(4)} виглядав нас на порозі, спокійно курячи
свою \VUgras{люльку}.

%%K t07-c1-03-2 p
Він запросив нас увійти, і мене не довелося просити двічі, щоб я всівся перед
вогнем із виноградної лози, що тріщав у вогнищі. Тоді мені став дарма різкий
\VUgras{північний вітер}, від якого рипіла стара \VUgras{вивіска}
\textsuperscript{(3)} і який ніс чорними клубами \VUgras{дим}
\textsuperscript{(1)}, що виходив із димаря.

%%K t07-c1-04-1 p
4. --- Була година обіду. Не гаючись, ми віддали належне простій, але смачній
їжі, яку нам подали (*).

%%K t07-tit-3 sub
\VUpk{10.2pt}{40}{Кілька відвідин.}
\VUsaut{3.0mm}

%%K t07-c1-05-1 p
5. --- Після обіду я пішов з батьком, страховим агентом, до його клієнтів.
Спершу ми відвідали \VUgras{коваля} \textsuperscript{(5)}, який живе поруч із
корчмою. Він підковував коня. Я подивувався силі його підручного, який міцно
тримав \VUgras{копито} коня коло свого шкіряного фартуха, поки коваль важкими
ударами молота прибивав \VUgras{підкову} \textsuperscript{(6)}.

%%K t07-c1-06-1 p
6. --- Потім ми пішли до сільського \VUgras{шевця} \textsuperscript{(7)}, старого
Якова. Хоч вивіскою йому правив чудовий \VUgras{чобіт}, він удовольнявся тим, що
підшивав стару пару стоптаних черевиків.

%%K t07-c1-06-2 p
Старий зі сміхом сказав нам: «У місті я був „чоботарем“, і в мене не було
замовників. Тут я „латальник“, і роботи хоч відбавляй».

%%K t07-c1-07-1 p
7. --- Він говорив нам про свого сусіда \VUgras{голяра} \textsuperscript{(8)},
клієнти якого всі невдоволені. Його \VUgras{бритви} завжди щерблені, а
\VUgras{ножиці} ніколи до пуття не ріжуть. Але що він єдиний голяр у селі,
людям, звісно, нічого не лишається, як голитися і стригтися в нього. До того ж
він людина скупа і непривітна.

%%K t07-c1-08-1 p
8. --- На щастя, інші \VUgras{сусіди} на нього не схожі. \VUgras{Стельмах}
\textsuperscript{(9)}, приміром, людина добра, роботяща й послужлива. Віддати
йому воза в направу --- сама втіха. Робота в нього завжди чисто зроблена.

%%K t07-c1-09-1 p
9. --- Не нарікав старий Яків і на \VUgras{лимаря} \textsuperscript{(10)}, якому
інколи помагає шити \VUgras{збрую}, коли робота не терпить.

%%K t07-c1-09-2 p
Батько спитав його про родину. «Усі здорові. Онука моя в \VUgras{школі}
\textsuperscript{(12)}. \VUgras{Учителька} \textsuperscript{(13)} нею дуже
вдоволена; вона добра \VUgras{учениця} \textsuperscript{(14)}. Не те що цей
шибеник, мій син; у нього тільки ігри на думці. \VUgras{Учитель}
\textsuperscript{(15)} ніяк не змусить його вивчити бодай щось».

%%K t07-tit-4 sub
\VUsaut{3.0mm}
\VUpk{10.2pt}{40}{Сільські балачки.}
\VUsaut{3.0mm}

%%K t07-c1-10-1 p
10. --- Тоді старий розповів нам сільські новини. Збиралися будувати нову школу,
бо та, що в \VUgras{сільській управі}, стала вже надто тісна. Це було б
неважко, бо досить було купити частину саду \VUgras{мірошника}. Бідоласі це було
б дуже до речі. Від самих холодів \VUgras{жорна} \VUgras{млина}
\textsuperscript{(16)} не могли молоти зерно, бо \VUgras{колесо}
\textsuperscript{(17)} схопило \VUgras{льодом} \textsuperscript{(25)} на
\VUgras{струмку} \textsuperscript{(23)}.

%%K t07-c1-11-1 p
11. --- «До речі про будівлі: чи бачили ви нашу нову \VUgras{церкву}
\textsuperscript{(18)}? Вона дуже мальовнича під сніговою шапкою. \VUgras{Ворони}
\textsuperscript{(19)}, які не впізнають більше своєї старої дзвіниці, понуро
сідають на велике дерево \VUgras{цвинтаря} \textsuperscript{(20)}».

%%K t07-c1-12-1 p
12. --- Згадка про цвинтар нагадала шевцеві про тих батькових знайомих, що
померли за минулий рік. «Скільки \VUgras{могил} \textsuperscript{(21)}, добродію
мій, прибуло до давніших під \VUgras{кипарисом} \textsuperscript{(22)}! Смерть
узяла нашого старого нотаря. Усе село було на його \VUgras{похороні}».

%%K t07-c1-13-1 p
13. --- «Он його наступник ковзається на струмку. Він недавно заходив, щоб я
полагодив ремінець його \VUgras{ковзана} \textsuperscript{(26)}, який
луснув».

%%K t07-c1-14-1 p
14. --- «А он там, на березі струмка, новий \VUgras{польовий сторож}
\textsuperscript{(27)}. Він колишній жандарм і пострах сільських шибеників. Вони
розбігаються, ледве вгледять його \VUgras{гетри} \textsuperscript{(29)} і його
\VUgras{кашкет} \textsuperscript{(28)}».

%%K t07-c1-15-1 p
15. --- «А! он старий Йосип везе \VUgras{віз хмизу} \textsuperscript{(30)}
пекареві. --- І справді, сказав батько, впізнаю його запряг \VUgras{мулів}
\textsuperscript{(31)}. Мені треба з ним переговорити, і я з цього скористаюся.
Прощавай, старий Якове».

%%K t07-c1-16-1 p
16. --- Якраз коли ми виходили, сільські діти вибігали зі школи й кидалися на
\VUgras{льодову доріжку} \textsuperscript{(33)}, важачи впасти й зламати щось при
\VUgras{падінні} \textsuperscript{(34)}. Один із них узяв у стельмаха свої
\VUgras{санчата} \textsuperscript{(32)}, посадив у них товариша і побіг.

%%K t07-c1-17-1 p
17. --- Інші кинулися до \VUgras{замету} \textsuperscript{(37)} біля
\VUgras{мосту} \textsuperscript{(40)} і взялися обстрілювати \VUgras{снігову
бабу} \textsuperscript{(39)}, яку зліпили напередодні і якій дали капелюха,
люльку і шкільну торбу через плече.

%%K t07-c1-18-1 p
18. --- Їм був дарма крижаний вітер і холоднеча. У більшості не було навіть
\VUgras{шарфа} \textsuperscript{(35)}, а щоб зігрітися після гри в
\VUgras{сніжки} \textsuperscript{(38)}, їм досить було жмені гарячих
\VUgras{каштанів} \textsuperscript{(42)}, куплених у старої Жакліни,
\VUgras{перекупки}, яка дрижить від холоду під своєю надто тонкою
\VUgras{шаллю} \textsuperscript{(43)}.

%%K t07-tit-5 sub
\VUsaut{3.0mm}
\VUcentre{12.0pt}{0}{\textit{Друга сцена.}}
\VUsaut{3.0mm}

%%K t07-tit-6 sub
\VUpk{10.2pt}{40}{Ферма навесні.}
\VUsaut{4.0mm}

%%K t07-c2-01-1 p
1. --- При всіх радощах зими ми з глибокою втіхою стрічаємо повернення
\VUgras{весни}.

%%K t07-c2-01-2 p
Яка втішна картина --- двір ферми ввечері, у травні!

%%K t07-c2-02-1 p
2. --- На обрії поволі сідає \VUgras{сонце} \textsuperscript{(1)}, заливаючи
\VUgras{поля} \textsuperscript{(3)} багряними \VUgras{променями}
\textsuperscript{(2)}. Роботящий \VUgras{орач} \textsuperscript{(6)} поспішає
зорати свій \VUgras{лан} \textsuperscript{(4)}.

%%K t07-c2-02-2 p
Своїм \VUgras{плугом} \textsuperscript{(5)}, якого тягне міцний \VUgras{кінь}
\textsuperscript{(7)}, він проводить \VUgras{борозну} \textsuperscript{(8)}, у
яку \VUgras{сівач} \textsuperscript{(9)} жменями розкидає \VUgras{насіння}, що
дасть золоті колоски.

%%K t07-c2-03-1 p
3. --- \VUgras{Косарі} \textsuperscript{(11)} зі своїми косами скосили лугову
траву, гребці висушили \VUgras{сіно} \textsuperscript{(10)}, перевертаючи його
\VUgras{вилами} \textsuperscript{(12)} і граблями, і ось вони вертаються додому.

%%K t07-c2-03-2 p
Одні йдуть попереду важкого воза, тягненого волами, несучи знаряддя на плечі.
Інші, ліниві, вигідно вмостилися на запашному сіні.

%%K t07-c2-04-1 p
4. --- Проходячи, вони приязно вітаються з \VUgras{садівником}
\textsuperscript{(16)}, який працює в \VUgras{саду} \textsuperscript{(13)},
підрізаючи \VUgras{плодові дерева} і прищеплюючи \VUgras{груші}
\textsuperscript{(14)}. \VUgras{Сливи} \textsuperscript{(15)} у цвіту, а на добре
угноєному \VUgras{городі} \textsuperscript{(17)} городина, капуста і салата вже
гарні.

%%K t07-c2-04-2 p
На низькій огорожі гарний \VUgras{павич} \textsuperscript{(18)} розпускає хвіст,
щоб показати своє чудове пір'я.

%%K t07-tit-7 sub
\VUpk{10.2pt}{40}{Двір ферми.}
\VUsaut{3.0mm}

%%K t07-c2-05-1 p
5. --- Двері \VUgras{стайні} \textsuperscript{(19)} відчинені, і видно ясла й
\VUgras{підстилку} \textsuperscript{(23)} \VUgras{кобили} \textsuperscript{(21)},
яку \VUgras{конюх} \textsuperscript{(20)} щойно розпряг. \VUgras{Лоша}
\textsuperscript{(22)}, таке кумедне на довгих ногах, іде за матір'ю, пустуючи.

%%K t07-c2-06-1 p
6. --- Коло \VUgras{криниці} \textsuperscript{(29)} \VUgras{корови}
\textsuperscript{(25)} йдуть пити з дерев'яного \VUgras{корита}
\textsuperscript{(26)}, що править їм за водопій. \VUgras{Теля}
\textsuperscript{(24)} користається з нагоди посмоктати, а \VUgras{віл}
\textsuperscript{(27)}, зачувши добрий запах корму, весело мукає і гордо
підіймає голову з могутніми \VUgras{рогами} \textsuperscript{(28)}.

%%K t07-c2-07-1 p
7. --- Усі при роботі. \VUgras{Робітниця} \textsuperscript{(32)} тримає
\VUgras{мотузку} \textsuperscript{(31)} криниці. Вона черпає воду \VUgras{відром}
\textsuperscript{(30)}, щоб напоїти худобу. Коло \VUgras{клуні}
\textsuperscript{(33)} чоловік згрібає розсипану солому, а перед дверима
\VUgras{ферми} \textsuperscript{(34)} стара \VUgras{пряха} \textsuperscript{(35)}
зі своєю прядкою і \VUgras{кужелем} пряде \VUgras{льон} або \VUgras{коноплі}, як
за давніх днів.

%%K t07-tit-8 sub
\VUsaut{3.0mm}
\VUpk{10.2pt}{40}{Фермер.}
\VUsaut{3.0mm}

%%K t07-c2-08-1 p
8. --- \VUgras{Фермер} \textsuperscript{(36)} керує всією цією роботою і дає
своїм людям вказівки. Він велить прибрати під повітку \VUgras{борону}
\textsuperscript{(40)} і \VUgras{кирку} \textsuperscript{(39)}, лишені коло
\VUgras{буди} \textsuperscript{(38)} \VUgras{собаки} \textsuperscript{(37)}.

%%K t07-c2-09-1 p
9. --- Його гуки полохають \VUgras{голуба} \textsuperscript{(41)}, який чимдуж
летить до \VUgras{голубника} \textsuperscript{(42)}, і \VUgras{курей}
\textsuperscript{(43)}, які поспішають назад до \VUgras{курника}
\textsuperscript{(44)}.

%%K t07-c2-10-1 p
10. --- Перед \VUgras{кошарою} \textsuperscript{(48)} старий \VUgras{пастух}
\textsuperscript{(45)} приганяє свою \VUgras{отару} \textsuperscript{(49)}.
Тримаючи свою \VUgras{ґирлиґу} \textsuperscript{(46)}, яка не розлучається з ним,
як і вицвіла \VUgras{блуза} \textsuperscript{(47)}, він веде в загороду
\VUgras{баранів} \textsuperscript{(50)}, \VUgras{овець} \textsuperscript{(53)},
\VUgras{ягнят} \textsuperscript{(54)}, \VUgras{кіз} \textsuperscript{(51)} і
\VUgras{козенят} \textsuperscript{(52)}. Його вірний \VUgras{пес}
\textsuperscript{(55)} Аргус іде останній і гавкотом підганяє відсталих.

%%K t07-c2-11-1 p
11. --- Увесь цей гамір і метушня не тривожать спокою доброї \VUgras{дійної
корови} \textsuperscript{(56)}, яка подзенькує своїм \VUgras{дзвіночком}
\textsuperscript{(57)}, поки її доять перед дверима \VUgras{корівника}
\textsuperscript{(58)}.

%%K t07-c2-12-1 p
12. --- Вона немов розуміє, що має дати молока, щоб \VUgras{фермерка}
\textsuperscript{(60)} могла збити \VUgras{масло} в \VUgras{маслобійці}
\textsuperscript{(61)} або зробити добірні \VUgras{сири} \textsuperscript{(64)},
які обсихають на дошці, покладеній упоперек \VUgras{діжки}
\textsuperscript{(63)}.

%%K t07-c2-13-1 p
13. --- Усю \VUgras{молочарню} \textsuperscript{(59)} тримає в чистоті
\VUgras{наймит} \textsuperscript{(65)}, який щойно вимив \VUgras{бідони}
\textsuperscript{(62)} у великому відрі, що несе.

%%K t07-tit-9 sub
\VUsaut{3.0mm}
\VUpk{10.2pt}{40}{Пташиний двір.}
\VUsaut{3.0mm}

%%K t07-c2-14-1 p
14. --- У своїх важких дерев'яних черевиках він ходить доволі незграбно і тому
полохає \VUgras{індика} \textsuperscript{(68)}, \VUgras{качок}
\textsuperscript{(66)}, \VUgras{качурів} \textsuperscript{(67)} і \VUgras{гусей}
\textsuperscript{(69)}, які широко роззявляють \VUgras{дзьоби}
\textsuperscript{(70)} і зчиняють тривожний ґвалт.

%%K t07-c2-14-2 p
\VUgras{Півень} \textsuperscript{(71)}, войовничіший, підіймає свій червоний
\VUgras{гребінь} \textsuperscript{(72)}, струшує пір'ям \VUgras{хвоста}
\textsuperscript{(73)} і видає дзвінке «кукуріку».

%%K t07-c2-15-1 p
15. --- \VUgras{Квочка} \textsuperscript{(74)} порпається в \VUgras{купі гною}
\textsuperscript{(81)} зі своїми \VUgras{курчатами} \textsuperscript{(75)}.
\VUgras{Порося} \textsuperscript{(78)} біжить до матері, величезної
\VUgras{свині} \textsuperscript{(77)}, яку видно коло хліва.

%%K t07-c2-16-1 p
16. --- У своїх клітках \VUgras{кролі} \textsuperscript{(79)} жадібно їдять ніжну
\VUgras{траву} \textsuperscript{(80)}, яку приносить їм фермерова дочка. Женя
дуже любить своїх кролів і щодня, взявши в курнику свіжі \VUgras{яйця}
\textsuperscript{(76)}, приходить навідати своїх маленьких друзів.
\VUsaut{6.0mm}
\VUfilet{22mm}
